\subsection{边界正则性}\label{sec:elliptic-boundary-regularity}

\PDEIndexMark{idx:boundary-regularity-nonsmooth-domain-warning-introduction}前一节说明, 若数据和系数足够光滑, 则椭圆问题的弱解在问题所在区域的内部具有 ``可以预期的光滑性\textquotesingle\textquotesingle. 以下例子表明, 若边界不够光滑, 则解未必光滑到问题所在区域的边界. (读者还应注意, 这不是某种奇怪的 ``人为构造\textquotesingle\textquotesingle\ 的反例; 所考察的区域只是有一个角点.)

\begin{Example}
\label{ex:nonsmooth-sector-boundary-regularity}
设 $\bm{i}$ 和 $\bm{j}$ 构成 $\mathbb{R}^2$ 的一组正交归一基, $(r,\theta)$ 为 $\mathbb{R}^2$ 的标准极坐标, 由映射
\begin{equation}\label{eq:corner-example-polar-coordinate-map}
\widehat{\bm{x}}(r,\theta):=r\bm{e}_1(\theta),
\end{equation}
定义, 其中
\begin{equation}\label{eq:corner-example-first-polar-basis-vector}
\bm{e}_1(\theta):=\cos\theta\bm{i}+\sin\theta\bm{j},
\end{equation}

\begin{equation}\label{eq:corner-example-second-polar-basis-vector}
\bm{e}_2(\theta):=-\sin\theta\bm{i}+\cos\theta\bm{j}.
\end{equation}
回顾对实值函数 $f(r,\theta)$, 可以如下计算梯度和 Laplacian:
\begin{equation}\label{eq:corner-example-polar-gradient-formula}
\nabla f=f_r\bm{e}_1+\frac{1}{r}f_\theta\bm{e}_2,
\end{equation}

\begin{equation}\label{eq:corner-example-polar-laplacian-formula}
\Delta f=\frac{1}{r}\frac{\partial}{\partial r}\left(r\frac{\partial f}{\partial r}\right)+\frac{1}{r^2}\frac{\partial^2f}{\partial\theta^2}.
\end{equation}
此外, Hessian 或二阶梯度矩阵为
\begin{equation}\label{eq:corner-example-polar-hessian-matrix}
\begin{pmatrix}
f_{rr}&\dfrac{1}{r}f_{r\theta}-\dfrac{1}{r^2}f_\theta\\
\dfrac{1}{r}f_{r\theta}-\dfrac{1}{r^2}f_\theta&\dfrac{1}{r}f_r+\dfrac{1}{r^2}f_{\theta\theta}
\end{pmatrix}.
\end{equation}
现在考虑 Laplace 方程具有非齐次边界条件的如下问题. 设 $0<\beta<2\pi$, 定义
\begin{equation}\label{eq:corner-example-planar-sector-domain}
\Omega_\beta:=\{\bm{x}\in\mathbb{R}^2\mid\bm{x}=\widehat{\bm{x}}(r,\theta),\quad0<r<1,\quad0<\theta<\beta\}.
\end{equation}
寻求 $u:\Omega_\beta\to\mathbb{R}$ 满足
\begin{equation}\label{eq:corner-example-harmonic-sector-equation}
\Delta u=0\quad\text{in }\Omega_\beta,
\end{equation}
及边界条件\[\begin{aligned}
u(r,0)&=0,&&0<r<1\\
u(r,\beta)&=0,&&0<r<1\\
u(1,\theta)&=\sin\frac{\pi\theta}{\beta},&&0<\theta<\beta.
\end{aligned}\]利用分离变量法, 可求得问题的解为
\begin{equation}\label{eq:corner-example-explicit-harmonic-sector-solution}
u(r,\theta)=r^{\pi/\beta}\sin\frac{\pi\theta}{\beta}.
\end{equation}
如前一节的结果所保证的, 对紧含于 $\Omega_\beta$ 的任意 $\Omega^{\prime}$, 这个函数都属于 $C^\infty(\Omega^{\prime})$. 然而, 注意若 $\beta>\pi$, 则解在原点附近像 $r^\alpha$ 那样衰减, 其中 $\alpha<1$. 为了解这一点的影响, 计算 $u$ 的梯度
\begin{equation}\label{eq:corner-example-explicit-solution-gradient}
\nabla u=\frac{\pi}{\beta}r^{\pi/\beta-1}\left(\sin\frac{\pi\theta}{\beta}\bm{e}_1(\theta)+\cos\frac{\pi\theta}{\beta}\bm{e}_2(\theta)\right)
\end{equation}
及其 $L^2(\Omega_\beta)$ 范数
\begin{equation}\label{eq:corner-example-finite-gradient-l2-norm}
\int_{\Omega_\beta}|\nabla u|^2\,\mathrm{d}\bm{x}=C\int_0^\beta\int_0^1r^{2(\pi/\beta-1)}r\,\mathrm{d}r\,\mathrm{d}\theta<\infty.
\end{equation}
可见, 如弱解的基本存在性定理所蕴含的, 解 $u$ 属于 $H^1(\Omega_\beta)$. 然而, 计算二阶梯度及其范数后可知, 若 $\beta>\pi$, 则 $u$ 不属于 $H^2(\Omega_\beta)$. 因此, 尽管具备前一节结果所保证的全部内部正则性, 却没有直到边界的正则性. 原因正在于边界缺乏光滑性.
\end{Example}
